\section{The continuous companion fibre and its margins}\label{sec:continuous}

The finite audits have a continuous companion in which the common-action obstruction is a linear-programming fact and its certificates are dual vectors. The instances below (Figure~\ref{fig:continuous}) are the smallest that exhibit the crossover, the saddle, the nonconvex gap and the Helly bound; the general theory, including the benchmark law of which Proposition~\ref{prop:bench} is the demand-$2$ instance, is developed in \cite{Abaee2026minimax}, and the continuous-to-finite bridge in \cite{Abaee2026cert}.

\begin{proposition}[benchmark fibre: crossover, margin law, radius]\label{prop:bench}
Two instruments $u_1,u_2\ge 0$ with caps $u_i\le c_i(Y)$, $c_1(Y)=\tfr32-\tfrac{Y-2}{10}$ and $c_2(Y)=\tfrac{59}{50}-\tfrac{Y-2}{10}$, must meet the demand $u_1+u_2\ge 2$; $Y$ is a pressure parameter. The problem is feasible if and only if $Y\le Y^{\ast}=\tfrac{27}{5}$. For $Y>Y^{\ast}$ the normalised dual vector $\lambda=(\tfr12,\tfr12)$ certifies infeasibility with margin
\[
\mu(Y)=1-\tfr12\bigl(c_1(Y)+c_2(Y)\bigr)=\frac{Y-27/5}{10}>0,
\]
e.g.\ at $Y=6$ the caps are $\tfrac{11}{10}$ and $\tfrac{39}{50}$, the cap sum $\tfrac{47}{25}$ and the margin $\tfrac{3}{50}$; at $Y=5$ the point $(\tfr65,\tfr45)$ is a feasibility witness. The margin is a robustness radius: raising both caps by $\delta$ restores feasibility at $Y$ if and only if $\delta\ge\mu(Y)$, so at $Y=6$ every per-cap increase below $\tfrac{3}{50}$ leaves the obstruction in place and the increase $\tfrac{3}{50}$ removes it.
\end{proposition}

\begin{proof}
Feasibility is $c_1(Y)+c_2(Y)\ge 2$, and $c_1+c_2=\tfrac{67}{25}-\tfrac{Y-2}{5}$ equals $2$ at $Y=\tfrac{27}{5}$. For $Y>Y^{\ast}$, $\lambda\cdot u\le\tfr12(c_1+c_2)<1=\lambda\cdot(2,0)$ for every admissible $u$, so no admissible $u$ meets the demand; the identity for $\mu$ is the linear cap law, and the radius statement is the same inequality with $c_i+\delta$.
\end{proof}

\begin{proposition}[static duality on a two-floor instance]\label{prop:duality}
Let $r_1(u)=\tfr25-u$ and $r_2(u)=u-\tfr35$ be two floor rows on $U=[0,1]$, and define the primal value $\max_{u\in U}\min_j r_j(u)$ and the dual value $\min_{\lambda\in\Delta}\max_{u\in U}\sum_j\lambda_j r_j(u)$.
\begin{enumerate}[label=(\alph*),leftmargin=2em,itemsep=1pt]
\item Both values equal $-\tfrac{1}{10}$, attained at the saddle $u=\tfr12$, $\lambda=(\tfr12,\tfr12)$; the negative value is the obstruction certificate: no $u$ satisfies both floors, and $\lambda$ prices the shortfall.
\item Convexity is load-bearing: on $U=\{0,1\}$ with payoff rows $(-1,+1)$ and $(+1,-1)$ the primal value is $-1$ and the dual value $0$.
\item A feasible contrast: with $r_2$ replaced by $u-\tfr15$ the two values are $+\tfrac1{10}$, attained at $u=\tfrac{3}{10}$; a certificate exists exactly where the value is negative, and a zero value is boundary-critical with no positive margin.
\end{enumerate}
\end{proposition}

\begin{proof}
(a) $\min_j r_j(u)=\min(\tfr25-u,u-\tfr35)$ is maximised where the rows cross, $u=\tfr12$, with value $-\tfrac1{10}$; for $\lambda=(\tfr12,\tfr12)$ the combination is the constant $-\tfrac1{10}$, which bounds the dual value above, and weak duality bounds it below. (b) Each column has a $-1$, so the primal is $-1$; $\lambda=(\tfr12,\tfr12)$ gives the constant $0$ and no $\lambda$ does better since the rows sum to $0$. (c) The rows cross at $u=\tfrac3{10}$ with value $\tfrac1{10}$ and the same $\lambda$ gives the constant $\tfrac1{10}$. Weak duality holds unconditionally; equality in (a) and (c) is von Neumann's theorem for affine rows on a convex compact $U$ \cite{vonNeumann1928}, and (b) is its failure without convexity. All values are verified on the rational grids of step $\tfrac1{100}$.
\end{proof}

\begin{proposition}[sparse witnesses and quantifier order]\label{prop:helly}
\begin{enumerate}[label=(\alph*),leftmargin=2em,itemsep=1pt]
\item An infeasible finite family of affine constraints on a convex domain of $\R^m$ has an infeasible subfamily of at most $m+1$ members (Helly \cite{Helly1923}), and $m+1$ is tight: $\{u_i\ge 1\}_{i\le m}$ together with $\sum_i u_i\le m-1$ is infeasible while every $m$ of its members are feasible. On the plane the triple $\{u_1\ge 1,\ u_2\ge 1,\ u_1+u_2\le 1\}$ is pairwise feasible (at $(1,1)$, $(1,0)$, $(0,1)$) and jointly infeasible.
\item Floors must be met by one instrument: with $f_1(u)=u-\tfr12$ and $f_2(u)=-u-\tfr12$ on $[-1,1]$, each floor is satisfiable ($u=1$, $u=-1$) but no $u$ satisfies both; $\forall j\,\exists u$ holds and $\exists u\,\forall j$ fails, with $\max_u\min_j f_j(u)=-\tfr12$.
\end{enumerate}
\end{proposition}

\begin{proof}
(a) is Helly's theorem applied to the half-spaces intersected with the domain; omitting $\sum u_i\le m-1$ take all $u_i=1$, omitting $u_j\ge 1$ take $u_j=0$ and the others $1$; the full family requires $m\le m-1$. (b) is a direct check; the value is attained at $u=0$.
\end{proof}

\begin{proposition}[certainty-equivalent drift]\label{prop:drift}
Let the stock evolve by the nonlinear law $g(s)=s^2$ applied to a certainty-equivalent lead of $0.1$. The naive drift $g(s+\tfr1{10})-g(s)=\tfrac{s}{5}+\tfrac1{100}$ is positive at every audited level ($\tfrac{21}{100},\tfrac{13}{50},\tfrac{31}{100},\tfrac{9}{25},\tfrac{41}{100}$ at $s=1,\tfr54,\tfr32,\tfr74,2$); the corrected law that subtracts $2s\cdot\tfr1{10}+\tfrac1{100}$ has zero drift at every level; and a correction error $e$ leaves the drift $2se+e^2$, whose sign is the sign of $e$ for $|e|<2s$. A certainty-equivalent floor is therefore exact only with the quadratic correction, and the direction of its failure is the direction of the correction error.
\end{proposition}

\begin{proof}
Expand $(s+\tfr1{10})^2$; the statements are polynomial identities, checked exactly.
\end{proof}

\paragraph{Margins as radii.}
The three continuous certificates come with margins that are robustness radii rather than slack numbers: the benchmark obstruction at $Y=6$ survives every per-cap perturbation below $\tfrac3{50}$ and not beyond; the timing identity of Corollary~\ref{cor:grid} has zero slack at the boundary cells $(2.0,1)$, $(3.0,2)$, $(4.0,3)$, where any increase of the damage rate flips the verdict, and positive slack $z_0-1-D_T$ elsewhere; and the drift law has the sign law of Proposition~\ref{prop:drift}. The finite kernels have an analogous but discrete robustness: the inclusions of Theorem~\ref{thm:mono} are stable under every perturbation, while each strictness rests on a named witness ($\st{12|21}$, $\st{12|31}$, $\st{13|21}$, $\st{13|31}$) whose status changes with a single transition.

\begin{figure}[t]
\centering
\includegraphics[width=\textwidth]{fig5_continuous.pdf}
\caption{The continuous fibre. Left: the cap sum $c_1(Y)+c_2(Y)$ against the demand $2$; the crossover $Y^{\ast}=27/5$ and the dual margin $\mu(Y)=(Y-27/5)/10$ beyond it. Right: the two floor rows of Proposition~\ref{prop:duality} and their lower envelope; the saddle at $u=1/2$ has value $-1/10$, and the feasible contrast (dotted) has value $+1/10$ at $u=3/10$.}
\label{fig:continuous}
\end{figure}
